\section{Empirical Strategy}
\label{sec:model}

\subsection{Pass-through from within-day surcharge discontinuities}
Let $T_{dm}$ be the mean credit-card tip on trips picked up on date $d$
in minute-of-day $m$, $S_{dm}$ the mean metered surcharge, and $B_{dm}$
the mean pre-tip total \emph{net} of the surcharge---the metered fare and
other components. Around a clock cutoff $c$ where a surcharge switches,
with running variable $r = m - c$ and $D = \mathbf{1}\{r \ge 0\}$, we
estimate
\begin{equation}\label{eq:rd}
  T_{dm} \;=\; \alpha + \pi D + \gamma B_{dm} + f(r) + \varepsilon_{dm},
\end{equation}
with $f(r)$ local linear on each side, cells weighted by card-trip
counts, and standard errors clustered by date. Two scalings of $\pi$
appear throughout. Dividing by the \emph{statutory} dose gives the
intention-to-treat pass-through per legislated dollar, the paper's
headline estimand. Instrumenting the measured surcharge $S_{dm}$ with
$D$ gives $\rho^{IV} = \pi/\phi$, with $\phi$ the first-stage step; with
a continuous treatment this is an average causal response
\parencite{angrist1995two}, and it carries an exclusion restriction
discussed in Section~\ref{sec:results-read}.

The preferred specification differences the discontinuity of
equation~\eqref{eq:rd} against the weekend one at the same clock, where
no surcharge changes \parencite{grembi2016difference}. With $W_d$ an
indicator for weekday dates and slopes interacted with $W_d$,
\begin{equation}\label{eq:didisc}
  T_{dm} \;=\; \alpha + \pi D + \omega W_d + \delta\,(D \times W_d)
    + \gamma B_{dm} + f(r) \times W_d + \varepsilon_{dm},
\end{equation}
and $\delta$---the difference in discontinuities---divided by the
statutory dose is the headline number. Differencing removes any 4:00pm
discontinuity common to weekdays and weekends.

\subsection{Why the non-surcharge base must be controlled}
The term $\gamma B_{dm}$ is not cosmetic; it is the difference between a
credible estimate and a biased one. At 4:00pm the tip base rises by
\$2.49, of which \$1.95 is the surcharge. The remaining \$0.53 is the
non-surcharge base---chiefly the metered fare, which jumps \$0.49
because rush-hour trips are slower and dearer. Omitting $B_{dm}$
attributes those fare-driven tip increases to the surcharge and inflates
the estimate by roughly a third. Two diagnostics confirm the correction.
First, the weekend placebo---the same clock on days when no rush
surcharge exists---falls from \$0.088 to \$0.024 once $B_{dm}$ is
included, which is what a correctly specified design should do to a
placebo. Second, the corrected estimates
agree across designs that have no reason to share a bias, including one
identified purely in space (Section~\ref{sec:results-fee}). Under the
weekday--weekend difference the control is nearly redundant: dropping it
moves the estimate by about a cent (Section~\ref{sec:results}).

\subsection{Identifying assumption}
For the single difference, identification requires rider composition
and tipping norms to evolve smoothly through the clock cutoff
conditional on the fare. For the preferred difference-in-
discontinuities the requirement is weaker: whatever jumps at 4:00pm for
reasons other than the surcharge---shift changes, light, the rhythm of
the working day---must jump equally on weekends, where the clock
carries no surcharge. The residual threat is a weekday-specific,
non-surcharge discontinuity at exactly the statutory minute; placebo
cutoffs at surcharge-free minutes and a permutation over candidate
cutoffs bound it. The assumption is also checkable, and
Section~\ref{sec:results-6am} shows where it fails: at 6:00am on
weekdays the dawn compositional shift is discontinuous in precisely the
prohibited way, and the design falls apart there while its weekend twin
does not. We take the failure as evidence the diagnostics bind, not as
grounds to discard them.

\subsection{Dose, direction, and transportability}
The December 2022 increase supplies a second dose at each clock, so
comparing discontinuities across regimes tests whether one $\rho$
governs both. The within-day structure also signs the shocks: surcharges
switch \emph{on} at 4:00pm (weekdays) and 8:00pm (weekends), and the net
surcharge \emph{falls} at 8:00pm on weekdays when the rush surcharge
ends and the smaller overnight surcharge begins. Comparing upward with
downward changes tests whether pass-through is symmetric.

Finally, the 2019 New York State congestion surcharge and the January
2025 MTA congestion toll moved the same tip base at boundaries in space.
Because the fee columns exist only after each policy begins, treated
areas must be defined geographically: we learn the set of pickup zones
where each fee is actually charged from post-period incidence and apply
those zone sets to the pre-period, then estimate
difference-in-differences on zone-by-fare-bin cells with zone-fare and
month fixed effects, clustering by pickup zone. This design shares no
variation with the clock designs, so agreement between them is a
genuine out-of-sample test.
